% ============================================================
% PAGE 1 -- TITLE PAGE
% ============================================================
\begin{titlepage}
\begin{center}

\vspace*{0.5cm}

\includegraphics[width=3.2cm]{vit-ap-university-logo.png}

\vspace{0.3cm}

{\Large\bfseries
VIT-AP UNIVERSITY
}

\vspace{0.3cm}

{\large
Amaravati, Andhra Pradesh, India
}

\vspace{0.8cm}

{\Large\bfseries
COMPILER DESIGN LAB REPORT
}

\vspace{0.4cm}

{\large
\textbf{Subject Name: Compiler Design Lab}
}

\vspace{0.15cm}

{\large
\textbf{Course Code: CSE3024}
}

\vspace{1cm}

{\LARGE\bfseries
Machine Learning-Based Loop Optimization\\
for High-Performance Compiler Design
}

\vspace{1.2cm}

\textbf{Submitted by}

\vspace{0.4cm}

\begin{tabular}{ll}
\textbf{Poojith} & \textbf{Regd. No.:} 24BCB7244 \\[0.4cm]
\textbf{Student 2 Name} & \textbf{Regd. No.:} 24BCB7161 \\
\end{tabular}

\vspace{1.2cm}

\textbf{Under the Guidance of}

\vspace{0.3cm}

\textbf{Dr. Prakash Kumar Sarangi}

\vspace{1cm}

{\large
School of Computer Science and Engineering (SCOPE)\\
VIT-AP University\\
Amaravati, Andhra Pradesh, India
}

\vspace{0.4cm}

{\large
Academic Year 2026--2027
}

\end{center}
\end{titlepage}

% ============================================================
% PAGE 2 -- ABSTRACT
% ============================================================
\section*{Abstract}

High-performance computing programs spend most of their time in loop nests, and
the speed of the generated code depends on optimisation choices such as
vectorisation, unrolling and the use of the target instruction set. Production
compilers make these choices with fixed, hand-written heuristics that cannot
adapt to the individual loop. This report designs and evaluates an
ML-based loop-optimisation decision model that sits in front of an unmodified
GCC compiler.

A compiler front end is implemented from scratch: a lexical analyser, a
recursive-descent parser, a semantic and loop analyser, and a three-address-code
generator. From the analysed loop nest the front end computes \numFeatures\
static features, such as trip counts, exact memory strides, loop-carried
dependence distance, reductions, arithmetic intensity and working-set size. A
supervised model maps the features to one of \numConfigs\ GCC optimisation
configurations. The model is trained on \numKernels\ C loop kernels from
\numFamilies\ HPC-style families (vector updates, reductions, recurrences,
1-D/2-D stencils, matrix multiplication, 2-D traversals, branchy and strided
loops). Every kernel was compiled with every configuration (\numCompiles\
compilations) and executed on the target machine, and the label is the best
configuration that preserves the numerical result.

The evaluation uses the speed-up of the generated code. With random 5-fold
cross-validation the best model (\selModel) reaches a geometric-mean speed-up of
\selGMCV$\times$ over \texttt{-O2}, which is \selFracCV\% of the gain of the oracle
(\oracleGM$\times$). This must be compared honestly with static choices: always using
\texttt{-O3} gives \gmOThree$\times$, always using \texttt{-O3 -march=native} gives
\gmONative$\times$, and always using \texttt{-O3 -march=native -ffast-math} gives
\gmOFast$\times$. On this machine the learned model is therefore only marginally better
than the most aggressive static configuration in cross-validation, and on loop families
never seen in training (leave-one-family-out) it is worse (\selGMLOFO$\times$), which shows the
limits of a small corpus. The model is more useful when the binary must be portable and
IEEE-safe: among the four non-native, non-fast-math configurations it reaches \pCV$\times$
(oracle \pOracle$\times$) against \pOThree$\times$ for \texttt{-O3}, which causes \pOThreeSlow\
slow-downs of more than 5\%. Feature extraction and prediction take about \featMs\ ms and
\predMs\ ms per loop, small compared with a GCC compilation (\gccCompileMs\ ms).

\textbf{Keywords:}
Compiler Design, Loop Optimization, Lexical Analysis, Syntax Analysis,
Dependence Analysis, Machine Learning, High-Performance Computing, GCC.


\newpage

% ============================================================
% PAGE 3-4 -- INTRODUCTION
% ============================================================
\section{Introduction}

A compiler is a system software component that translates a program written in
a source language into an equivalent program in a target language. For
high-performance computing (HPC) workloads such as dense linear algebra,
stencil codes, signal processing and scientific simulation, almost all of the
execution time is spent inside a small number of \emph{loop nests}. The quality
of the machine code generated for these loops therefore determines the
performance of the whole application \cite{allen2002,bacon1994}.

Modern optimising compilers contain many loop transformations: loop unrolling,
loop interchange, loop fusion and fission, tiling, loop-invariant code motion,
strength reduction and automatic vectorisation \cite{aho2006,cooper2007,muchnick1997}.
Whether a transformation pays off depends on the loop's trip count, memory
access pattern, data type, data-dependence structure, working-set size and on
the target micro-architecture. Production compilers decide with hand-written
heuristics and fixed cost models that were tuned once, on a set of benchmarks, for
an average machine. These heuristics are cheap but can be far from the best choice
for an individual loop, and tuning them by hand is slow and has to be repeated
for every new processor \cite{wang2018,ashouri2018}.

Machine learning offers an alternative. Instead of writing the heuristic, we
\emph{learn} it: a number of loops are compiled and timed under different
optimisation choices, and a model is trained to predict the best choice from
static features of the loop that the compiler front end can already compute
\cite{stephenson2003,stephenson2005,fursin2011}. The learned decision model
replaces the hand-written profitability test of the optimiser.

This report presents a complete, working implementation of this idea for the
Compiler Design Lab. A small compiler front end (lexical analyser, recursive
descent parser, semantic and loop analyser, three-address-code generator) is
written from scratch and used to extract \numFeatures\ numeric features from C
loop nests. A supervised model then selects one of \numConfigs\ GCC optimisation
configurations for the loop. The model is trained and evaluated on
\numKernels\ loop kernels that were actually compiled and executed on the
target machine.

\subsection{Motivation}

Compiler design is usually taught as a fixed pipeline of phases. The motivation
of this project is to show that the \emph{outputs} of the early phases (tokens,
syntax tree, symbol table, intermediate code) are exactly the information that
an optimisation decision needs, and that the last phase, optimisation, can be
driven by a data-driven model instead of an ad-hoc rule. A practical
understanding of this connection links compiler theory with the current
research direction of ML-guided compilation used in industrial toolchains
\cite{trofin2021,chen2018}.

The topic is also relevant because HPC users routinely tune compiler flags by
trial and error: compile with \texttt{-O3}, then with \texttt{-march=native},
then with unrolling, and measure. A model that predicts a good configuration
directly from the source of a loop can save this search time.

\subsection{Problem Statement}

Given the source code of a loop nest, the compiler must decide which
optimisation configuration to apply. This report addresses the following
problem:

\begin{quote}
\textit{Given a C loop nest $L$, can a decision model that uses only static
features computed by a compiler front end predict which of $K$ candidate
optimisation configurations gives the lowest execution time, so that the
compiled code is faster than with the default configuration (\texttt{-O2})
and close to the best possible (oracle) configuration?}
\end{quote}

Formally, let $C=\{c_1,\dots,c_K\}$ be the configurations and $T(L,c)$ the
measured execution time of loop $L$ compiled with $c$. We want a function
$f:\mathbb{R}^{d}\rightarrow C$ on the feature vector $\phi(L)\in\mathbb{R}^{d}$ with
\begin{equation}
f(\phi(L)) \approx \arg\min_{c\in C} T(L,c).
\end{equation}

\subsection{Objectives}

\begin{enumerate}
    \item To implement the front-end phases of a compiler (lexical, syntax and
    semantic analysis, intermediate code generation) for a C loop-nest subset.

    \item To perform loop analysis (trip counts, memory strides, loop-carried
    dependences, reductions) and convert the result into a numeric feature
    vector.

    \item To build a data set of loop kernels labelled with the best GCC
    optimisation configuration, obtained by real compilation and timing.

    \item To train and compare ML decision models, and to evaluate them by the
    speed-up of the generated code and not only by classification accuracy.

    \item To package the result as a command-line tool that advises the
    compiler flags for a new loop.
\end{enumerate}

\subsection{Contributions}

The major contributions of this report are as follows:

\begin{itemize}
    \item \textbf{Compiler front end for loop analysis:}
    A lexer, recursive-descent parser, semantic analyser and TAC generator for
    a C loop subset, implemented from scratch in Python
    (Section~\ref{sec:method}).

    \item \textbf{Loop feature set:}
    \numFeatures\ static features derived from the AST and TAC, including
    exact memory strides computed from array declarations and
    dependence distances computed from affine subscripts.

    \item \textbf{Measured data set:}
    \numKernels\ kernels in \numFamilies\ families, each compiled with
    \numConfigs\ configurations and timed with an interleaved, noise-aware
    protocol (\numCompiles\ compilations).

    \item \textbf{Performance evaluation of the decision model:}
    Evaluation by realised speed-up with two protocols, random cross-validation
    and leave-one-family-out, against static and oracle baselines
    (Section~\ref{sec:results}).
\end{itemize}

\newpage

% ============================================================
% PAGE 5-7 -- RELATED WORK
% ============================================================
\section{Related Work}

Compiler design has been extensively studied since the development of the
first automatic programming systems. Aho et al.\ presented comprehensive
foundations for lexical analysis, parsing, semantic analysis, intermediate
code generation, optimization, and code generation \cite{aho2006}. Cooper and
Torczon described optimization techniques and their implementation in modern
compiler infrastructures \cite{cooper2007}. Appel discussed compiler
construction and code generation techniques \cite{appel2004}.

\subsection{Traditional Compiler Models}

Traditional compiler architectures use a sequence of analysis and
synthesis phases. Lexical analysis converts character sequences into tokens.
Syntax analysis organizes tokens into parse trees. Semantic analysis verifies
contextual correctness. Intermediate representations are then optimized before
target-code generation. LR and LL parsing techniques have been extensively used
for context-free grammars; the front end in this report uses a recursive-descent
(LL) parser with precedence climbing for expressions.

\subsection{Compiler Optimization and Loop Transformations}

Optimization attempts to improve program performance while preserving
program semantics. Classical techniques include constant folding, constant
propagation, common subexpression elimination, dead-code elimination, loop
optimization, strength reduction, and control-flow optimization
\cite{muchnick1997}. For loops, the central transformations are
\emph{unrolling} (replicating the body to reduce loop overhead and expose
instruction-level parallelism), \emph{interchange} and \emph{tiling}
(re-ordering iterations to improve cache locality \cite{wolf1991}) and
\emph{vectorisation} (executing several iterations with one SIMD instruction).
Their legality is decided by data-dependence analysis, which computes
dependence distances between array references \cite{allen2002,bacon1994}.
Polyhedral frameworks automate the search of legal loop transformations
\cite{bondhugula2008}.

Machine-independent optimisation operates on intermediate representations
without depending on a particular processor architecture. Machine-dependent
optimisation considers target-specific characteristics such as registers,
instruction sets (e.g.\ AVX2/AVX-512), cache behaviour and pipeline constraints.

\subsection{Machine Learning in Compiler Optimization}

Wang and O'Boyle survey the use of machine learning in compilers and conclude that
learning replaces hand-built heuristics for decisions such as unroll factors,
inlining, phase ordering and device mapping \cite{wang2018}; Ashouri et al.\ survey
compiler auto-tuning \cite{ashouri2018}. Stephenson and Amarasinghe used supervised
classification to predict loop unroll factors from loop features
\cite{stephenson2005}, and Stephenson et al.\ evolved compiler heuristics with
genetic programming \cite{stephenson2003}. Agakov et al.\ used learned models to
focus iterative optimisation on promising transformations \cite{agakov2006}, and
Milepost GCC integrated feature extraction and machine learning into GCC itself
\cite{fursin2011}. Cummins et al.\ learned optimisation heuristics directly from
source text with deep networks \cite{cummins2017}. MLGO brings
learned policies into LLVM \cite{trofin2021}, and TVM uses a learned cost model
to tune tensor kernels \cite{chen2018}.

Most of this work either modifies the compiler internals or needs large benchmark suites. The approach in
this report is deliberately lightweight: the decision model sits \emph{in front of}
an unmodified compiler (GCC) and chooses among its existing flag sets, which makes it
easy to reproduce in a laboratory setting.

\subsection{Compiler Technology for Machine Learning}

Machine learning frameworks generate computational graphs containing
operations such as matrix multiplication, convolution and activation functions.
Compiler-based approaches such as XLA and MLIR optimise these graphs before
execution \cite{abadi2016,lattner2021}. This is the opposite direction of the
present work (compilers \emph{for} ML versus ML \emph{in} compilers), but both
rely on the same intermediate-representation and loop-optimisation machinery.

\subsection{Survey of Related Approaches}

\begin{table}[H]
\centering
\caption{Survey of Compiler Optimisation Decision Approaches}
\label{tab:survey}
\small
\begin{tabularx}{\textwidth}{p{2.6cm}p{2.5cm}X p{2.9cm}}
\toprule
\textbf{Approach} & \textbf{Decision} & \textbf{Major Contribution} &
\textbf{Limitation}\\
\midrule

GCC / LLVM \texttt{-O} levels
& All passes
& Mature, fixed heuristics and cost models
& Tuned for an average program and machine\\

Polyhedral (Pluto)
& Loop tiling, parallelism
& Automatic legal transformation search for affine loops
& Restricted to affine loop nests\\

Unroll-factor prediction \cite{stephenson2005}
& Unroll factor
& Supervised classification from loop features
& Single transformation, internal compiler change\\

Iterative compilation + ML \cite{agakov2006}
& Transformation sequence
& Models focus the search on good sequences
& Needs many compile-run cycles\\

Milepost GCC \cite{fursin2011}
& Flags / passes
& Static features and learning inside GCC
& Large infrastructure\\

DeepTune \cite{cummins2017}
& Device mapping, coarsening
& Learns features from raw source
& Needs large training sets\\

MLGO \cite{trofin2021}
& Inlining, register allocation
& Learned policies in production LLVM
& Tied to the compiler version\\

\textbf{This work}
& Flag configuration of loops
& Own front end + features + model in front of unmodified GCC
& Small corpus, one machine\\

\bottomrule
\end{tabularx}
\end{table}

\subsection{Limitations of Existing Approaches}

Although learned and heuristic optimisers provide significant performance
benefits, several limitations remain:

\begin{enumerate}
    \item Hand-written heuristics are tuned for an average machine and do not
    adapt to the loop or to the processor.
    \item Machine-specific optimisations (e.g.\ \texttt{-march=native}) reduce
    portability of the generated binary.
    \item Some transformations change floating-point results
    (e.g.\ \texttt{-ffast-math}), so a decision model must be checked for
    numerical safety.
    \item Learned models can mis-predict on loop patterns that were absent from
    the training data; this must be tested explicitly.
    \item Timing labels are noisy on shared machines, so labelling has to be
    robust to measurement noise.
    \item Optimisation increases compilation effort, and the decision model
    itself must be cheap compared with compilation.
\end{enumerate}

The proposed methodology addresses items 3--6 explicitly: it rejects
configurations that change the numerical result, labels with a noise
tolerance, evaluates on unseen loop families, and measures the overhead of
the model.

\newpage
